%!TeX root = ./../MusterAbschlussarbeit.tex

%##########################################################
% Inhalt
%##########################################################

\clearpage
\chapter{Auswertung und Ausblick}


\section{Bewertung der Ergebnisse}
\section{Schritte zur Verbesserung des Agenten}
\section{Diskussion}

Dieser Abschnitt stellt den Schlusspunkt der Arbeit dar. In diesem Abschnitt (und im Diskussionteil) erwartet der Leser, 
dass er Antworten auf die in der Einleitung formulierten Fragestellungen findet und sich vergewissert, 
dass diese wirksam verteidigt wurden und mit der von Ihnen formulierten \enquote{These} übereinstimmen.

Die Hauptziele der Diskussion bestehen darin, eine Analyse Ihrer gesammelten Ergebnisse zu präsentieren, 
Ihre Ergebnisse angemessen darzustellen und eine Einschätzung der Bedeutsamkeit Ihrer Arbeit zu geben. Beachten Sie hier,
den Unterschied zur Diskussion im vorherigen Kapitel. Diskutieren Sie hier vor allem den Wert und die Bedeutung Ihrer Ergebnisse, auf Basis der 
Interpretationen aus dem vorherigen Kapitel. Beziehen Sie sich gern auch auf Ihr beschriebenes Problem und Ihr Ziel.

Jede wichtige Schlussfolgerung, die Sie im \enquote{Ergebnisteil} gezogen haben, 
muss hier erneut behandelt werden. Eine gewisse Anzahl von Wiederholungen ist unvermeidlich. Darüberhinaus, 
die Ergebnisse anderer Forschungsarbeiten, müssen mit eindeutigen Verweisen auf auffindbare Literaturquellen versehen sein.

Der Fazit-Teil kann als eine kurze Zusammenfassung Ihrer Diskussion betrachtet werden. 
Der Leser muss sich hier schnell einen Überblick über den Inhalt und die Bedeutung der Arbeit als Ganzes verschaffen.

Dieser Abschnitt soll einen Überblick präsentieren und dient dazu, dem Hauptteil Ihrer Arbeit den letzten Schliff zu geben. 
Die Schlussfolgerung kann auch Hinweise auf ein mögliches zukünftiges Werk enthalten.

\underline{Aber \textbf{mindestens} folgende Punkte}\\
Inhalt des fünften Kapitels (im Allgemeinen):
\begin{itemize}
    \item Reflektieren Sie hier nun die Ergebnisse der Tests aus dem letzten Kapitel
    \item Ordnen Sie diese in den Gesamtkontext ein... gut/schlecht? Was kann man verbessern/anders machen? Schätzen Sie auch ab was Veränderungen bringen könnten.
    \item Was wurde durch die Ergebnisse gezeigt? Geben Sie eine bewertende (selbstkritische) Aussage ab zu Ihrem Schaffen
    \item Geben Sie einen Ausblick was nun folgen sollte/könnte.
\end{itemize}



\lstinputlisting[caption=C-Code Beispiel,label={lst:C-Code-Beispiel},language=C]{Programmcode/test.c}
\lstinputlisting[caption=Python-Code Beispiel,label={lst:Python-Code-Beispiel},language=Python]{Programmcode/test.py}

\textbf{Es ist dabei auch darauf zu achten, dass Quelltexte wenn möglich natürlich nicht über das Seitenende gehen sollten, wie z.B. beim Python-Code Bsp. Setzen Sie daher auch manuell
Seitenumbrüche.} Quelltexte welche eine Seite übersteigen, sind ohnehin eher als Anhang zu betrachten $\rightarrow$ siehe dazu Beschreibung im Anhang.



\lstinputlisting[caption=HTML-Code Beispiel,label={lst:HTML-Code-Beispiel},language=HTML]{Programmcode/test.html}

Nachfolgend noch ein einfaches Beispiel ein Bild einzubinden. Da es sich um Bild\textit{unterschriften} handelt, gehört diese somit \underline{unter} die Abbildung, anders als bei Tabellen\textit{überschriften}.\\
Im \LaTeX-Quelltext sieht man die Optionen für das Einbinden des Bildes. Neben \texttt{scale} wäre auch \texttt{width} möglich mit dem Parameter \texttt{linewidth} oder \texttt{textwidth}.\\
Achten Sie vor allem auf die Lesbarkeit der auf dem Bild befindlichen Informationen, unter der stetigen Annahme, dass es sich um ein ausgedrucktes Dokument handelt. Dort gibt es keinen Zoom.

Empfehlenswert wäre soweit möglich mit skalierbaren Vektorgrafiken zu arbeiten. Allerdings müssten Sie diese vor dem Aufruf von LaTeX in PDFs wandeln (inkscape IM-logo.svg -o IM-logo.pdf) oder Inkscape installieren und LaTeX mit --shell-escape aufrufen.

\begin{figure}[!t]
    \centering
    \includegraphics[scale=0.6]{Bilder/logos/IM-logo}
    \caption{Das Logo von Informatik und Medien }
\end{figure}

\begin{figure}[!t]
    \centering
    \includesvg[width=0.6\columnwidth]{Bilder/logos/IM-logo.svg}
    \caption{Das Logo von Informatik und Medien }
\end{figure}

Bedenken Sie bei allen Beschriftungen, egal ob Abbildungen, Tabellen oder Quelltexte, dass Sie, sofern diese nicht Ihrer eigenen Schaffenskraft entsprangen,
diese referenzieren. Im Beispiel der Tabelle \ref{tab:Beispieltabelle} sieht man, dass sich der Beschreibungstext direkt an der Tabelle von dem im Tabellenverzeichnis unterscheidet,
genau um den Punkt der Referenz.
